% ============================================================
%  Poul Martin Møller, "What Is Philosophy?"
%  English translation of the Introduction ("Indledning", ms. leaves
%  1r--10v) to Møller's lectures on the history of ancient philosophy,
%  University of Copenhagen, c. 1833--35.
%
%  SOURCE TEXT: Møller's own manuscript, in the edition of Det Danske
%  Sprog- og Litteraturselskab, Poul Martin Møllers Skrifter
%  (http://pmm.dsl.dk/forelaesninger/txt096.1.xml; text ed. Finn Gredal
%  Jensen; version 1.0, 2011-10-10). The Danish reading text is
%  generated from DSL's TEI by ../../../pmm2tex.py into
%  ../pmm/hvad-er-philosophie.tex (kept local, gitignored, until DSL
%  agrees to a mirror):
%    python3 pmm2tex.py "texts/moller/pmm/Den ældre Philosophies Historie I.xml" \
%      texts/moller/pmm/hvad-er-philosophie.tex \
%      --until "Om den græske Philosophies Oprindelse"
%  Final readings of the manuscript only; Møller's deletions are not
%  translated. Square brackets in the text are Møller's own.
%
%  This is NOT the text of Efterladte Skrifter (vol. 3, 1843; vol. 4,
%  3rd ed. 1856, pp. 3--17), where F. C. Olsen supplemented Møller's
%  draft from students' lecture notes. The partial transcription.tex in
%  this folder (from Børge Madsen's selection) follows that printed text.
%
%  \opage{1r} etc. = manuscript leaf. "No 2" etc. at the head of a
%  sheet, and the heading Møller wrote at the top of each new sheet,
%  are kept as in the manuscript.
%
%  Places where the manuscript (as DSL gives it) is unfinished or odd,
%  translated as the sense requires and marked with a comment:
%  "hvad er Penge." (1v); "give enfuldkommen" (4v); "er er gaaet" (6r);
%  "opbevare Mærkeligste" (7r, words wanting); "er fuldkommen, han
%  udretter" (8v, sentence unfinished); "have have" (9r).
% ============================================================
\documentclass[12pt]{article}
\usepackage[utf8]{inputenc}
\usepackage[T1]{fontenc}
\usepackage{libertinus}
\usepackage{libertinust1math}
\usepackage{textalpha}
\usepackage[english]{babel}
\usepackage[protrusion=true,expansion=true]{microtype}
\usepackage[margin=1.4in]{geometry}
\usepackage{setspace}
\usepackage{fancyhdr}
\usepackage{marginnote}
\usepackage[colorlinks=true,linkcolor=black,urlcolor=black,
  pdftitle={What Is Philosophy?},
  pdfauthor={Poul Martin Møller}]{hyperref}
\setlength{\headheight}{15pt}
\pagestyle{fancy}
\fancyhf{}
\fancyhead[LE,RO]{\thepage}
\fancyhead[RE]{\textit{What Is Philosophy?}}
\fancyhead[LO]{\textit{Poul Martin Møller}}
\setlength{\parindent}{1.5em}
\setlength{\parskip}{0pt}
\setstretch{1.3}
\newcommand{\opage}[1]{\marginnote{\footnotesize\textit{[#1]}}}

\title{\textbf{What Is Philosophy?}\\[0.5em]
  \large Introduction to the Lectures on the History of Ancient Philosophy\\
  University of Copenhagen, c.\ 1833--35}
\author{Poul Martin Møller}
\date{Translated from the Danish of Møller's manuscript, as edited by\\
  Det Danske Sprog- og Litteraturselskab (\url{http://pmm.dsl.dk})}

\begin{document}
\maketitle

\begin{flushright}\opage{1r}No.\end{flushright}

\begin{center}[Introduction.]\end{center}

\noindent What is philosophy? --- One might expect this question to be answered by anyone who means to lecture on the history of philosophy. --- It is easy to give a good many equally correct definitions of philosophy, but they would serve only to remind someone who knows what philosophy is of what he thinks by it. If he did not already have a concept or a representation of philosophy, he would not, on hearing a single proposition, come all at once to recognize what philosophy is. In an ordinary definition of a concept one usually gives a superordinate concept (\textit{dist.\ gen.}) and a characteristic mark (\textit{dist.\ specif}), by which the concept is distinguished from other, coordinate concepts. I would thus have to determine philosophy with the help of a couple of concepts that stand in close connection with it; but how should someone who thought nothing by the word philosophy be able to think anything thorough on hearing the names of the related concepts? No one can have these in a thorough way unless he has philosophized, and so knows what philosophy is. Suppose I were to say: philosophy is a science of the ultimate grounds of things, or of the eternal; still no one who did not know at all what philosophy is could know clearly what I mean by the ultimate grounds of things\opage{1v} or the eternal.

To this is added yet another difficulty. Every proposition in which one would express the concept of philosophy is so one-sided that anyone who knows how to think something by it is at once led to change his expression and to put forward opposite determinations. I say philosophy is the science of the eternal; but since it is not to treat of the abstract eternal but of the eternal that reveals itself, I can say with just as much truth that it is the science of the finite. --- I say that it is the science of the ultimate grounds of things; but it is just as much the science of the consequences of these grounds, since grounds first show themselves as grounds in their consequences. Thus I am led further on by every one-sided predicate, and can only be stopped arbitrarily at a one-sided definition. The complete expression of the concept of philosophy is the presentation of philosophy itself, and every explanation that stops before this whole presentation is accomplished is properly only partly correct.

So not much is gained by the usual, supposedly exact and scientific manner of putting a definition first. This formality is therefore skipped, and it is presupposed --- as of course it can be with good reason --- that everyone who hears these presentations knows what philosophy is, and knows it just as well as after hearing any definition whatever. What is money. % ms.: "hvad er Penge." -- a jotted comparison, apparently: one knows what money is without a definition.

On the other hand, it might perhaps be fitting to say something about the difference between philosophy and the doctrine of religion on the one side, and on the other side the scientific\opage{2r} thinking that is not philosophical. But I dare do this only in a couple of words, since the distinctions to be made here have so often been made in recent times that to anyone who has merely read a fair number of German and Danish reviews and prefaces to philosophical treatises they must soon seem like old, familiar rigmaroles.

In the first place, philosophy differs from a doctrine of religion that is not philosophical. They do not differ in their content, since God and his relation to the world are the chief objects of both. But the doctrine of religion may differ from philosophy in its form, and may do so more or less. The more the content of the religious consciousness is presented merely in a symbolic way, the more it puts images of fancy and significant myths in place of pure determinations of thought, the more it differs from philosophy. The rational being who receives with faith pictorial presentations of the nature of the Godhead and ingenious tales of the life of the gods and the origin of the world receives something more or less foreign, something he cannot find in the same way within himself, because what is offered him are not universally valid productions of the rational thinking that is the same in all. Sensuous representations of the eternal do not belong to philosophy, however profound a world-view may underlie them. This is remarked because people have often wanted to draw the history of such representations into the history of philosophy, especially where there has been talk\opage{2v} of the philosophy of the Orient. Philosophy has to do only with thoughts, and even if the most profound thoughts lie hidden in myths and images of fancy, one is here outside the territory of philosophy; for there the thoughts are not to lie hidden but to be definitely pronounced. On the other hand, ingenious myths and religious representations can prepare the way for philosophy and serve as incitements to philosophical thinking, and history teaches that they often have. The doctrine of religion comes nearer to philosophy when it contains thoughts in the form of thoughts; but when the religious propositions are merely accepted in faith and in trust in the infallibility of him who handed them down, then this kind of doctrine too, however true and profound it may otherwise be, cannot be counted as a philosophical presentation of the truth. For that it is required that thinking itself be the only source of knowledge to which appeal is made. The presentation of the truth must be universally valid in order to be philosophical, so that every thinking individual who follows the presentation can follow it with full conviction, because it is an expression of the train of thought that is necessarily determined by universal reason. Philosophy, then, differs from the doctrine of religion that appeals to another authority than reason itself, whether that doctrine moves within the precincts of representation or of thoughts.

\begin{flushright}\opage{3r}No.\ 2.\end{flushright}

\begin{center}\emph{Continuation of the Introduction to the Hist.\ of Philos.}\end{center}

Philosophy must further be distinguished from the scientific thinking that is not philosophical. This does not share its content with philosophy, but it has this in common with philosophy, that it brings self-active thinking to bear and regards it as a way to the knowledge of truth. It also has in common with philosophy that it goes beyond what is given in sensation and strives to know the universal that underlies the phenomena, brings them under universal concepts, and strives to know the laws that hold for their changes. But insofar as it is a thinking about something given, and aims at ordering and connecting a sum of representations of existing things found in the world of sense, it is not yet properly philosophical thinking. --- The scientific thinking that falls outside philosophy does indeed point toward the philosophical idea, since an obscure representation of the unity in the whole of existence guides [all its endeavors,] its whole activity; but it seeks to arrive at knowledge of this unity by comparing the finite and connecting it under more general surveys, or even by leaving out sensuous determinateness, abstracting from it. --- In its whole activity it makes a regulative use of thinking's own concepts, the a priori concepts or categories; to that extent it is partly by itself, or has to do with itself, besides the material given from without; but it has not come to consciousness of\opage{3v} the nature and essence of these concepts and of their relation to the nature and essence of things. From this it follows, partly, 1 that it uses these concepts at random, without heeding their relation to one another: a multitude of categories are used according to a recollection of how others use them; but these are regarded as fixed presuppositions, as the eyes of the understanding, with which one is to see, but which are not themselves to be seen. 2. Next, it follows that one does notice the consistent procedure of one's thinking and its presupposition of concepts that are not given by sensation, but takes them to be merely subjective and not valid for things. --- In philosophical thinking, on the other hand, thinking comes to consciousness that its nature and essence is the nature and essence of existence, and that by thinking merely according to its own laws it knows the eternal and true determinations of existence. It comes to consciousness that what must necessarily be thought also is as one is compelled to think it. The idea that served merely for regulative use in empirical thinking is then recognized to be the right point of departure, and the indeterminate starting point that one must determine more closely, by thinking's own authority, in order to arrive at a rational knowledge of things.

Thus we have, in a relative way, bounded the concept of philosophy on two sides, and excluded several products of thought from the spiritual content that we shall attend to in developing the history of philosophy.

\opage{4r}The history of philosophy naturally hangs together with the general history of mankind. Free thinking has in many ways been modified, hindered and furthered by the external circumstances under which it has developed. Whoever would fully explain to himself the shape it has taken among a given people at a given time therefore not seldom needs knowledge of that people's character and history. We shall develop this somewhat more closely:

Philosophy --- such as it exists, actually comes forth --- in a human consciousness always has a definite connection with the philosopher, with his whole spiritual life and circle of representations; but this in turn is, in the first place, largely determined by [the natural individuality and the stage of cultivation peculiar to the national stock to which he belongs, as also by] the natural character and historically developed condition of his national stock. The people's religious faith, its constitution, its scientific culture exert a necessary (i.e.\ unavoidable) influence on the direction his philosophy takes. In the second place, his own temperament and his empirical character come much into consideration. All these external factors one must know if one is to know the circle of his consciousness; and this in turn furnishes the context that must be known if a particular philosophy is to be fully grasped. A philosopher's words are not to be understood lexically. It therefore furthers clear insight into a given philosopher's full meaning to know the history and condition of his fatherland, and to know the events of his own life. The more life and clarity the presentation of a philosophy has, the more it bears the stamp of the particular circumstances under which it developed, and the deeper the roots it has struck in the life of its time: it has\opage{4v} partly stood in harmony with the opinions of its time, partly been directed polemically against them; and through knowledge of such circumstances much that is peculiar in a given philosophy first becomes properly explicable.

The history of philosophy is thus in a certain respect a fragment torn out of an organically connected whole, because it has a decided connection with the history of the whole human race. To be able to give a complete presentation of it, one must therefore study not merely the history of the sciences but universal history as such. The recognition that study of history is required to find one's way in the history of philosophy has brought several mere scholars, who had no interest in philosophical speculation, to take possession of this province [i.e.\ to take on this task, as a work belonging to their province]. Brucker. Tiedemann. Indeed, people have even gone so far as to maintain that whoever had no philosophical world-view of his own was best fitted to report the philosophers' doctrines as they are historically found, without falsifying them by additions of his own. This is a fundamentally false opinion, and no notice would have been taken of it here had it not been put forward even in the most recent times. Whoever would present the history of philosophy is not even able to sift his material out of the mass of literature unless he himself has philosophical cultivation and philosophical interest. (One sees this from the fact that in the history of philosophy one finds left out what is of a high degree of philosophical interest, and other things dragged in that wholly lack it (Tiedemann. Brucker.)

\begin{flushright}\opage{5r}No.\ 3\end{flushright}

\begin{center}Continuation of the Introduction to the Hist.\ of Philos.\end{center}

\noindent Even supposing that he knew from tradition which works were generally regarded as philosophical, he would still not be able, without philosophical insight, to separate what is essential in them from what is inessential, and still less, by an ingenious divination, to integrate the fragments that are left to us of certain philosophical systems.

Whoever would study the history of philosophy in a scientific way must then himself bring a philosophical spirit, and know the questions that the philosophical desire for knowledge demands be answered. Whoever makes the philosophical systems an object of indifferent apprehension, systems in which he himself has a merely historical interest, will in general take note of precisely what strikes his eye as strange and absurd: that is, of precisely what is unphilosophical in them, or is apprehended by him as such. The philosophical systems cannot be reduced mechanically to a smaller scale: one cannot set a copyist to make extracts of them as of a bundle of legal documents.\footnote{Those who would apprehend the history of philosophy without prejudices, and who believe they do so by not bringing any philosophy with them themselves [believe they secure themselves against them by taking care not to bring philosophy with them], nevertheless bring prejudices: namely, those of common human understanding.}

But it takes not only philosophical spirit and philosophical cultivation to find the thread in the history of philosophy: it takes many-sided receptivity and many-sided philosophical cultivation. Whoever regards all philosophy\opage{5v} from the one-sided standpoint of a restricted system sees everything in a skewed light and overlooks much. Friis --- Tennemann. % Friis: so in the ms.; J. F. Fries is meant.
If one now objects to this that no mortal is able to raise himself to the speculative Hlidskialf from which all philosophical systems could be surveyed so that one sees them in their true relation to one another, it by no means follows that the demand must be given up, since it lies in the nature of the case and cannot be refused. At most it would follow that no mortal can give a complete presentation of the history of philosophy. That is certainly true; but no one who works in the field and knows what the history of philosophy is can refrain from striving toward this goal, which reason has set.

The chief interest that the study of the history of philosophy has is that it teaches how philosophy came to be. It is thus the most thorough study of philosophy itself, because it lets us go through in thought all the periods of human reason that it has had to traverse in order to reach the stage of knowledge at which it now works. One learns from it that the philosophers' endeavors have not been fruitless; that the matter does not stand as is generally assumed by those who regard philosophy from an exoteric standpoint, that the earlier systems have vanished without trace or result to make room for those that followed; but that true philosophical inquiry is a work connected through many centuries, an unbroken chain\opage{6r} whose later links presuppose the earlier as their necessary condition. That a higher guidance has led mankind in its scientific striving is clearly seen in this: that the individuals did not know the building they were working on, and yet out of their united endeavors a whole and connected work has come forth. --- The history of philosophy has proceeded according to the unchangeable laws of rational thinking, along the same road that it would in the main have to travel anew if the human race now living were annihilated and a new race once more worked its way up from the state of nature to developed philosophical knowledge. It would have to travel this road anew because the road is grounded in the nature of rational thinking, as everyone comes to recognize by repeating it in his thoughts. For he will then find that every advance brought a new advance with it, that one one-sidedness was always corrected by another, until the third standpoint was found, on which the 2 conflicting views were seen in their relative truth. The thread that runs through the whole history of philosophy has never been broken; when for a time it seemed to have been dropped, inquiry has nonetheless always been resumed, and in such a way that it attached itself to the train of thought begun in an older time. When, for example, philosophical thinking began again in the Middle Ages,\opage{6v} it began with the race making itself familiar, by historical studies, with the knowledge antiquity had already acquired: Plato. Aristotle. But after the years of apprenticeship were over, the Christian age brought forth a new and deeper philosophy.

The history of philosophy thus shows, in grand style, the history of human consciousness, which the individual must itself go through on a smaller scale in order to appropriate, with complete comprehension, the knowledge of truth that belongs to the age in which it lives. % ms. "som det selv maa giennemgaae": det, presumably the individual.
 But insofar as the course of development that philosophy has had to take is the necessary course of human consciousness, it must also, in its main features, be determinable in advance, because reason is like itself in the individual as in the race. To that extent there can be talk of constructing the history of philosophy a priori, where the expression ``history'' is then taken in a less proper sense. If a history of philosophy is to answer to its name, inquiry must keep to what is factually given and confirmed by historical criticism. Insight into the necessary development that follows from the nature of thinking must serve only as an aid to finding that connection in mankind's philosophical thinking which can be historically proved by valid testimony. It is naturally a fault in the historian of philosophy if, to bring connection into his researches, he adds something out of his own head and his own fancy that what is given does not compel him to presuppose; and a still greater fault if, to defend the connection he believes he finds, he does violence to the historical facts.

\begin{flushright}\opage{7r}No.\ 4\end{flushright}

\begin{center}Continuation of the Introduction to the Hist.\ of Philosophy\end{center}

\noindent The fault spoken of here is thus the exact opposite of the one spoken of above, which consists in apprehending the philosophers' systems merely historically, as a heap of odd notions, dreams, cobwebs of the brain. If the history of philosophy is to have philosophical interest, the rational connection in it must be brought out: not merely the connection between the individual's particular doctrines, but the connection between the successive systems, and their relation to one another. But on the other hand the inquirer must not be overhasty and put his subjective train of thought in place of the development that history has handed down to us. --- Much that is handed down as the doctrines of philosophers undeniably hangs together only poorly with their principles, and much seems wholly irreconcilable with them. Nevertheless it has with good reason become a custom among historians of philosophy to preserve the most remarkable of what seems wholly arbitrary and accidental, % ms.: "opbevare Mærkeligste", words wanting.
especially where the reports of a famous man's teaching are few and fragmentary; for partly an explanation may perhaps be found for what many have found inexplicable, and partly we thereby gain a closer acquaintance with the philosopher's whole circle of representations, and that, as has already been remarked, gives us a clearer insight into the meaning of his presentation.

\footnote{\emph{This should perhaps stand elsewhere,} earlier.}But how far, in a complete presentation of the history of philosophy, should one carry this endeavor to make known the philosopher's whole circle of representations, in order to show the relation of his philosophical thoughts to his empirical consciousness? One easily sees that a middle way is to be kept here between\opage{7v} two extremes: on the one side one must take care not to lose sight of the philosophical interest and make the lecture a collection of historical notices and food for curiosity; on the other side one must not be too afraid of degrading philosophy by showing its connection with the whole personal life and with merely human opinions. This anxiety one sometimes sees Germans carry much too far, in that with a certain hauteur they would remove the concept entirely from the region of representation. In this their presentation differs by the whole breadth of heaven from the naive exposition of the Greeks.

We said above that the history of philosophy is the history of human consciousness, and we may on this occasion mention a view that has often been expressed and will be expressed more often still, since it has its ground in a property of human nature. Every independent and original thinker has, namely, the innate conviction that what he assumes on sufficient grounds must necessarily be right, whether others agree with him in it or not. It then lies near at hand for him to think and speak thus: ``Why should I search in old writings for what I have within myself? I bear within myself the source of true knowledge: the higher life of reason that moves in me shall lead my thoughts to a philosophical conviction that can in truth be called mine, since I myself have produced it by my own original power. At most I will go to school with one of the distinguished philosophers of the age, to learn on what\opage{8r} footing science stands; and then I will work only in my own free circle of thought, and not waste my time and cripple my strength of spirit by seeking from without what cannot be given from without. It is a proof of the impotence of the age, its lack of productivity, that its philosophical study has taken so decidedly historical a direction.''

So many have thought; and under certain circumstances it has been a great piece of good fortune that many have thought so and acted on this thought, for thereby many one-sided directions in philosophy have been pursued with force and consistency, and these are necessary preliminary labors for a many-sided presentation of philosophy. Every speculative genius can certainly, in a period of ferment, when many divergent philosophical systems are formed one after another, work out a system for himself that may be as true as many others. Such a period of ferment there has been once in Greece and once in modern times. Then, naturally, a one-sided idealism, for example, and a one-sided realism must appear, and the one such system is as right and as wrong as the other. But it is quite another matter when the ferment is over: then a philosophical day of judgment sets in, and the greatest geniuses then show their power in comprehending all these one-sided directions in their relative truth, and in producing a many-sided, organically connected system, in which the one-sided\opage{8v} systems rise again in a transfigured shape.
\par It goes in scientific life as in political: whoever, among savages who do not yet live in a rationally ordered society, makes his subjective conviction of what is rational and right prevail with force and success may become the founder of a state and be praised after his death as a hero; but whoever, in an already existing political order, would entirely tear down and overturn what prevails, without knowing the significance it has, in order to put in its place a constitution better in his own opinion, is entirely [\dots]; he accomplishes nothing rational. % ms.: "er fuldkommen, han udretter" -- the sentence is left unfinished.
So it goes in philosophy too: whoever, before the main principles of philosophy have been found, forcefully pronounces his philosophical knowledge as an autodidact from his one-sidedly subjective standpoint wins a place in the history of philosophy; but when the main principles of philosophy have been discovered, originality without knowledge of what the age knows has little significance: then even the greatest speculative genius comes to stand beneath his time, if he lacks either the knowledge or the receptivity to appropriate the knowledge his contemporaries have won. Herbart.

But the knowledge of the age in its whole many-sidedness can be attained only by a thorough study of the history of philosophy. It is the only road by which a philosophy developed in all directions can be comprehended, and it is the road the greatest philosophers of our time have taken. Thus Hegel, with extra-

\begin{flushright}\opage{9r}No.\ 5\end{flushright}

\begin{center}Introd.\ to the Hist.\ of Philos.\end{center}

\noindent ordinary genius and unusual learning, has striven to comprehend the eternal history of reason from the factual development of philosophy, and has carried out this plan, which others too have had, with a power with which no one else has carried it out. He has striven to fuse all philosophical systems into one system, in which they are all present as one organization: thus he has become the interpreter of the philosophers in the highest sense, in that he has comprehended their systems better than their inventors did, because he has apprehended them all from a higher standpoint. Such a historical philosophy is very different from an eclectic system, if one will use this oxymoron. The eclectic takes, by his subjective feeling, something from one man and something from another, without being able to comprehend it in its unity and connection.

But no presentation by a single man can take the place of the history of philosophy. Even if one were to credit some mortal with the many-sidedness to be an infallible judge over all the forms of philosophy, he still cannot, by his compressed presentation, give others a complete concept of all that he has seen into. In relation to the many world-views that are supposed to have their point of union in his system, his system is still at most a transfigured epitome, and as such abstract and general. For him, and for those who know the history of philosophy, the presentation can have a wealth of content that it does not have for one who\opage{9v} regards it torn loose from its relations. The greatest philosopher's work is a product of the philosophers' endeavors up to his time, but his work can be understood to the bottom only by one who knows the works of the preceding philosophers.

But can the history of philosophy then take the place of the presentation of philosophy? --- No, by no means. --- Philosophy is the knowledge that has been won through the historical development of philosophy, and without it the history of philosophy cannot be comprehended. Philosophy, looking back over its course, finds its \textit{membra disjecta} in history, and by philosophy they are brought into organic connection; without it one can find in history only a mass of conflicting elements. But philosophy unfolds itself in its many-sidedness by being found again, and finding itself again, in the manifold forms. The intension and the extension of knowledge stand in a direct relation to one another. The more moments in knowledge have stepped apart from one another, the more there lies in knowledge. Genuine philosophical knowledge can be regarded from two sides: from the standpoint of unity and from the standpoint of manifoldness. The system of philosophy is a unity that refers to a manifoldness, which is found in the history of philosophy. In presenting the system of philosophy one gives the unity, and in presenting its history the manifoldness. The unity is abstract for one who does not know the manifoldness, and the manifoldness is lost, uncomprehended, in contingency for one who does not know the system.

\opage{10r}But the objective history of philosophy is all the systems of which it consists, in their motley manifoldness and diversity of language. A presentation of the history of philosophy by a single man is in turn only a contribution to the study of the history of philosophy: it stands in the same relation to its objects as the description of nature to nature. It can be no more than an Adriadne's thread, by which one can be helped to find the way through the labyrinth that one must walk through oneself, in person (a miniature picture). % "Adriadne": so in the ms.

Thus we have seen what relation there is between philosophy and the history of philosophy, and what significance there is in the claim that only one who lacks productivity and originality can take an interest in the study of the history of philosophy. The philosophers who in modern times pass themselves off as original geniuses, and believe they deserve this predicate because they have not studied the history of philosophy, can in general be counted in 2 classes. To the first belong the really original geniuses, who stand beneath their time and pursue one or another one-sided direction in philosophy that mankind has already left behind. The others are original only in their own imagination: they are at bottom eclectics, who have, in various accidental ways --- by oral conversation, superficial reading, the study of reviews and prefaces --- come to know several\opage{10v} of the ideas of the age, and who now prepare from them a compound that they pass off as philosophy. The last class counts the most names. Neither Fichte nor Hegel nor Schelling announced himself in literature as an original: they became so gradually, without knowing it.

\begin{center}\rule{0.32\textwidth}{0.4pt}\end{center}

\end{document}
